% Options for packages loaded elsewhere
\PassOptionsToPackage{unicode}{hyperref}
\PassOptionsToPackage{hyphens}{url}
%
\documentclass[
]{article}
\usepackage{amsmath,amssymb}
\usepackage{iftex}
\ifPDFTeX
  \usepackage[T1]{fontenc}
  \usepackage[utf8]{inputenc}
  \usepackage{textcomp} % provide euro and other symbols
\else % if luatex or xetex
  \usepackage{unicode-math} % this also loads fontspec
  \defaultfontfeatures{Scale=MatchLowercase}
  \defaultfontfeatures[\rmfamily]{Ligatures=TeX,Scale=1}
\fi
\usepackage{lmodern}
\ifPDFTeX\else
  % xetex/luatex font selection
\fi
% Use upquote if available, for straight quotes in verbatim environments
\IfFileExists{upquote.sty}{\usepackage{upquote}}{}
\IfFileExists{microtype.sty}{% use microtype if available
  \usepackage[]{microtype}
  \UseMicrotypeSet[protrusion]{basicmath} % disable protrusion for tt fonts
}{}
\makeatletter
\@ifundefined{KOMAClassName}{% if non-KOMA class
  \IfFileExists{parskip.sty}{%
    \usepackage{parskip}
  }{% else
    \setlength{\parindent}{0pt}
    \setlength{\parskip}{6pt plus 2pt minus 1pt}}
}{% if KOMA class
  \KOMAoptions{parskip=half}}
\makeatother
\usepackage{xcolor}
\usepackage{longtable,booktabs,array}
\usepackage{calc} % for calculating minipage widths
% Correct order of tables after \paragraph or \subparagraph
\usepackage{etoolbox}
\makeatletter
\patchcmd\longtable{\par}{\if@noskipsec\mbox{}\fi\par}{}{}
\makeatother
% Allow footnotes in longtable head/foot
\IfFileExists{footnotehyper.sty}{\usepackage{footnotehyper}}{\usepackage{footnote}}
\makesavenoteenv{longtable}
\setlength{\emergencystretch}{3em} % prevent overfull lines
\providecommand{\tightlist}{%
  \setlength{\itemsep}{0pt}\setlength{\parskip}{0pt}}
\setcounter{secnumdepth}{-\maxdimen} % remove section numbering
% definitions for citeproc citations
\NewDocumentCommand\citeproctext{}{}
\NewDocumentCommand\citeproc{mm}{%
  \begingroup\def\citeproctext{#2}\cite{#1}\endgroup}
\makeatletter
 % allow citations to break across lines
 \let\@cite@ofmt\@firstofone
 % avoid brackets around text for \cite:
 \def\@biblabel#1{}
 \def\@cite#1#2{{#1\if@tempswa , #2\fi}}
\makeatother
\newlength{\cslhangindent}
\setlength{\cslhangindent}{1.5em}
\newlength{\csllabelwidth}
\setlength{\csllabelwidth}{3em}
\newenvironment{CSLReferences}[2] % #1 hanging-indent, #2 entry-spacing
 {\begin{list}{}{%
  \setlength{\itemindent}{0pt}
  \setlength{\leftmargin}{0pt}
  \setlength{\parsep}{0pt}
  % turn on hanging indent if param 1 is 1
  \ifodd #1
   \setlength{\leftmargin}{\cslhangindent}
   \setlength{\itemindent}{-1\cslhangindent}
  \fi
  % set entry spacing
  \setlength{\itemsep}{#2\baselineskip}}}
 {\end{list}}
\usepackage{calc}
\newcommand{\CSLBlock}[1]{\hfill\break\parbox[t]{\linewidth}{\strut\ignorespaces#1\strut}}
\newcommand{\CSLLeftMargin}[1]{\parbox[t]{\csllabelwidth}{\strut#1\strut}}
\newcommand{\CSLRightInline}[1]{\parbox[t]{\linewidth - \csllabelwidth}{\strut#1\strut}}
\newcommand{\CSLIndent}[1]{\hspace{\cslhangindent}#1}
\usepackage{bookmark}
\IfFileExists{xurl.sty}{\usepackage{xurl}}{} % add URL line breaks if available
\urlstyle{same}
\hypersetup{
  pdftitle={z4ai: Lossless, Random-Access Compression for Neural-Network Model Checkpoints},
  pdfauthor={Hyukjin Kwon},
  hidelinks,
  pdfcreator={LaTeX via pandoc}}

\title{z4ai: Lossless, Random-Access Compression for Neural-Network
Model Checkpoints}
\author{Hyukjin Kwon}
\date{June 2026}

\begin{document}
\maketitle
\begin{abstract}
The storage and distribution of neural-network checkpoints has become a
first-order systems cost, yet general-purpose compressors handle weight
tensors poorly and weight-specific codecs leave large, practically
important sources of redundancy untouched. We present z4ai, a lossless
codec for model checkpoints that operates on a single file and preserves
per-tensor random access. We first make explicit the
information-theoretic ceiling that binds \emph{every} lossless codec on
dense trained weights --- the exponent of a trained float carries only
\textasciitilde2.6 bits and the mantissa is near-random, capping ratios
at \textasciitilde1.51x for bf16 and \textasciitilde1.2x for fp32 ---
and show that z4ai ties the state of the art (ZipNN) there, as it must.
The contribution is not a new entropy coder but the integration of
whole-tensor long-distance matching, a lossless palette transform for
quantized weights shipped in wide float containers, and a
single-reference cross-checkpoint delta into one self-describing,
random-access container that is never worse than plain zstd. On the
redundancy regimes the entropy bound assumes away, this yields large
lossless gains over ZipNN: 4.72x vs 1.94x on INT8-in-fp32, 10.8x vs 4.6x
on INT4-in-fp32, 2.93x vs 1.67x on a realistic structured checkpoint,
and \textasciitilde30--184x on cross-checkpoint deltas. z4ai is
complementary to repository-scale storage systems such as ZipLLM,
targeting the pipelines and object stores that have no corpus-wide
deduplication service to build on.
\end{abstract}

\section{Introduction}\label{introduction}

Model registries, fine-tune families, and training runs produce large
numbers of checkpoints that are highly \emph{related} to one another,
and storing or transferring them at full size is increasingly expensive.
General-purpose compressors are a poor fit: in an IEEE-754 float the
low-entropy exponent and the high-entropy mantissa are interleaved
byte-by-byte, so a generic \texttt{zstd} pass over raw \texttt{fp32}
weights reaches only \textasciitilde1.06x. Weight-specific lossless
codecs --- ZipNN (Hershcovitch et al. 2024), NeuZip (Hao, Cao, and Mou
2024), DFloat11 (Zhang et al. 2025), DietGPU (Johnson 2022), and
Unweight (Cloudflare Research 2026) --- improve on this by separating
the float fields and entropy-coding the skewed exponent.

These codecs are, however, \emph{order-0} and operate within fixed-size
chunks, which leaves two practically large sources of redundancy
unaddressed. First, structure that spans a chunk boundary or a whole
tensor --- tied \texttt{embed\_tokens}/\texttt{lm\_head} matrices,
duplicated transformer blocks, multi-shard concatenations, and pruned
layers --- is invisible to a chunked, LZ-free coder. Second, the format
that models are \emph{actually shipped in}: quantized weights
(INT4/INT8/FP8 from GPTQ, AWQ, or \texttt{compressed-tensors}) are
routinely dequantized into a wide float container for deployment,
leaving only a small set of distinct values that a byte-grouping codec
cannot exploit. Finally, consecutive checkpoints in a training run
differ in only a few percent of their bytes, a redundancy exploited at
\emph{repository} scale by recent systems such as ZipLLM (Wang et al.
2026) but absent from drop-in codecs.

We present \textbf{z4ai}, a lossless codec that targets exactly these
regimes while remaining usable as a single-file library and command-line
tool with per-tensor random-access reads. Our contributions are:

\begin{enumerate}
\def\labelenumi{\arabic{enumi}.}
\tightlist
\item
  An explicit, reproducible statement of the \textbf{entropy ceiling} on
  dense trained weights, and a demonstration that z4ai ties ZipNN there
  within 0.3\% --- clarifying what no lossless codec can do.
\item
  The \textbf{integration} of whole-tensor matching, a lossless palette
  transform for dequantized quantized weights, and a single-reference
  cross-checkpoint delta into one self-describing container with a
  \emph{never-worse-than-zstd} guarantee.
\item
  An \textbf{evaluation} across dense, structured, sparse, quantized,
  real-checkpoint, and checkpoint-delta workloads, fully reproducible
  from the open-source repository.
\end{enumerate}

\section{Background: the entropy
ceiling}\label{background-the-entropy-ceiling}

A float is laid out as
\texttt{{[}sign\ \textbar{}\ exponent\ \textbar{}\ mantissa{]}}. In
trained weights the sign is approximately a fair coin (\textasciitilde1
bit), the mantissa is near-random (entropy close to its full width,
\textasciitilde7 bits for bf16), and only the exponent is strongly
skewed, with a measured entropy of roughly 2.57--2.74 bits against its
allocated 8 (Zhang et al. 2025; {``Lossless Compression of Neural
Network Components: Weights, Checkpoints, and {K/V} Caches in
Low-Precision Formats''} 2025). Entropy-coding the exponent and storing
the rest is therefore the whole game for a \emph{dense} tensor, and it
caps the achievable lossless ratio at approximately \textbf{1.51x for
bf16 and 1.2x for fp32}, regardless of the coder. ZipNN already sits at
that wall. This is an information-theoretic bound, not an engineering
gap: no lossless codec --- z4ai included --- can meaningfully exceed it
on dense, independently distributed weights. Larger ratios are only
available where the data carries redundancy the bound assumes away:
reduced precision, sparsity, repeated structure, or near-identical
neighbours in a checkpoint sequence. z4ai is designed entirely around
capturing that redundancy losslessly, and around \emph{not regressing}
on dense data where none exists.

\section{Design}\label{design}

z4ai decomposes each float tensor into sign, exponent, and mantissa
planes. The low-entropy exponent and sign planes are entropy-coded with
an interleaved range-ANS coder (Duda 2013; Giesen 2014), or with
\texttt{zstd} (Collet and Kucherawy 2021) when that is smaller for a
given stream; the noise-like mantissa is stored verbatim or
\texttt{zstd}-compressed. A \textbf{whole-tensor long-distance-matching}
pass (Ziv and Lempel 1977) then deduplicates repeated and tied weights
across the entire buffer --- the redundancy ZipNN's 256 KiB chunking
cannot see. For every stream, a \textbf{best-of selection} keeps the
smallest encoding, which guarantees the output is never larger than a
plain \texttt{zstd} pass.

Three transforms extend the codec beyond dense data:

\begin{itemize}
\tightlist
\item
  \textbf{Palette transform.} When a tensor in a wide float container
  holds only a small set of distinct values --- the signature of
  quantized weights dequantized for deployment --- z4ai relabels the
  values to a dense index plus a codebook. The transform is bijective
  and therefore lossless, and it recovers the original low bit-width
  that the wide container wasted.
\item
  \textbf{Sparsity path.} Pruned weights take a zero-aware encoding
  rather than paying to store runs of zero bytes.
\item
  \textbf{Cross-checkpoint delta.} \texttt{compress\_delta} /
  \texttt{model\_delta} implement name-aligned, per-tensor lossless
  deltas (copy / XOR-delta / full) against a reference checkpoint,
  robust to reordered or added tensors, so that checkpoint \emph{N}
  costs only the bytes that changed relative to \emph{N−1}.
\end{itemize}

The container is self-describing --- decoding is the exact inverse
driven entirely by the header, with no side-channel metadata --- and
carries a per-tensor index for random-access reads, storing tied tensors
once. An optional \texttt{effort="max"} tier adds a chunk-parallel
context-modeling backend (Alakuijala et al. 2018) that reaches
\emph{below} the order-0 floor on real transformer weights, falling back
to the fast default on genuinely incompressible data. General-purpose
float codecs such as ALP (Afroozeh, Kuffo, and Boncz 2024) and Pcodec
(Loncaric 2025) were evaluated but target smooth or decimal-origin
numeric data, not the full-entropy mantissa of trained weights, and were
not adopted.

\section{Evaluation}\label{evaluation}

All results below are measured on this repository's benchmark suite
(16-core machine, Python 3.14, \texttt{zstandard} 0.25.0, latest
\texttt{zipnn}, 32 MB per dtype, best-of-3 timing) and reproducible with
a single command per row. Every z4ai output is verified byte-exact.
Compression \emph{ratios} are load-independent and the primary metric;
throughput figures are reported for completeness and are sensitive to
machine load.

\subsection{Compression ratio}\label{compression-ratio}

\begin{longtable}[]{@{}llrrr@{}}
\toprule\noalign{}
Scenario & dtype & z4ai & ZipNN & zstd-3 \\
\midrule\noalign{}
\endhead
\bottomrule\noalign{}
\endlastfoot
Dense / i.i.d. & bf16 & 1.413 & \textbf{1.417} & 1.227 \\
Dense / i.i.d. & fp32 & 1.171 & \textbf{1.172} & 1.061 \\
Structured (repeated) & bf16 & \textbf{58.1} & 1.51 & 16.97 \\
Structured (repeated) & fp32 & \textbf{47.3} & 1.20 & 14.24 \\
Sparse (50\% zeros) & bf16 & \textbf{2.47} & 2.20 & 1.88 \\
Sparse (50\% zeros) & fp32 & \textbf{2.21} & 1.86 & 1.79 \\
Quantized INT8 → wide & bf16 & \textbf{2.39} & 2.07 & 1.79 \\
Quantized INT8 → wide & fp32 & \textbf{4.72} & 1.94 & 3.07 \\
Quantized INT4 → wide & bf16 & \textbf{5.41} & 3.87 & 3.91 \\
Quantized INT4 → wide & fp32 & \textbf{10.77} & 4.59 & 5.13 \\
\end{longtable}

On dense, i.i.d. weights z4ai ties ZipNN within 0.3\%, as the entropy
bound requires. The large gains appear exactly where redundancy exists.
(Caveat: on the quantized rows several ZipNN entries were \emph{not}
byte-exact on the tested build; the palette benchmark additionally
reports a stronger lossless \texttt{zstd-19} baseline, against which
z4ai's margin is more conservative than the table above.)

\subsection{Real and production
checkpoints}\label{real-and-production-checkpoints}

\begin{longtable}[]{@{}
  >{\raggedright\arraybackslash}p{(\linewidth - 6\tabcolsep) * \real{0.2143}}
  >{\raggedleft\arraybackslash}p{(\linewidth - 6\tabcolsep) * \real{0.2857}}
  >{\raggedleft\arraybackslash}p{(\linewidth - 6\tabcolsep) * \real{0.2857}}
  >{\raggedright\arraybackslash}p{(\linewidth - 6\tabcolsep) * \real{0.2143}}@{}}
\toprule\noalign{}
\begin{minipage}[b]{\linewidth}\raggedright
Workload
\end{minipage} & \begin{minipage}[b]{\linewidth}\raggedleft
z4ai
\end{minipage} & \begin{minipage}[b]{\linewidth}\raggedleft
ZipNN
\end{minipage} & \begin{minipage}[b]{\linewidth}\raggedright
Note
\end{minipage} \\
\midrule\noalign{}
\endhead
\bottomrule\noalign{}
\endlastfoot
bert-tiny, 17.7 MB fp32 (downloaded) & 1.188 & \textbf{1.202} & −1.2\%:
a single dense checkpoint is \textasciitilde i.i.d. \\
Production .safetensors, 201 MB bf16, tied embedding & \textbf{1.525} &
1.510 & +1.0\%: dedups the tied
\texttt{embed\_tokens}/\texttt{lm\_head} \\
Realistic full checkpoint, 107 MB bf16 (tied embeddings, shared blocks,
optimizer state, 50\%-pruned layer) & \textbf{2.93} & 1.67 & +75.7\%:
whole-buffer dedup of real structure \\
Checkpoint delta, bert-tiny bf16, 5\% of weights changed & \textbf{51.1}
& \textasciitilde1.7 & \textasciitilde30x smaller than from-scratch; 1\%
→ 184x, 20\% → 18x \\
\end{longtable}

\subsection{Throughput}\label{throughput}

\begin{longtable}[]{@{}lrr@{}}
\toprule\noalign{}
Codec & compress (MB/s) & decompress (MB/s) \\
\midrule\noalign{}
\endhead
\bottomrule\noalign{}
\endlastfoot
z4ai (i.i.d. bf16) & 1420 & 16700 \\
ZipNN & \textbf{8125} & \textbf{20020} \\
\end{longtable}

z4ai compresses several times slower than ZipNN's compiled-C core and
decompresses competitively --- a deliberate trade for the write-once,
read-many lifecycle of a stored checkpoint.

\section{Related work}\label{related-work}

\textbf{Weight-specific lossless codecs.} ZipNN (Hershcovitch et al.
2024) splits float fields and applies Huffman coding to the exponent in
256 KiB chunks; it is the closest dense baseline and, being chunked and
LZ-free, is exactly why it misses cross-tensor structure. NeuZip (Hao,
Cao, and Mou 2024) and DFloat11 (Zhang et al. 2025) entropy-code the
exponent with ANS and dynamic-length codes respectively, primarily for
memory-efficient training and GPU inference; DietGPU (Johnson 2022)
provides a GPU-native rANS float coder; Unweight (Cloudflare Research
2026) targets MLP weights for inference. All are order-0 on the float
fields and confirm the entropy decomposition of Section 2.

\textbf{Repository-scale storage.} ZipLLM (Wang et al. 2026) is the
closest \emph{system}: it deduplicates tensors across an entire model
corpus using content-defined chunking and model-family clustering, and
stores fine-tunes as lossless XOR deltas against their base, reporting
large whole-corpus savings. z4ai is \textbf{complementary}, not
competing: ZipLLM optimizes a managed repository as a service, whereas
z4ai is a self-contained per-file codec with random-access reads, an
explicit single-reference delta, and the palette transform --- for the
many pipelines and object stores that have no corpus-wide deduplication
service to build on. The broader landscape of low-precision and delta
compression for weights, checkpoints, and K/V caches is surveyed by
{``Lossless Compression of Neural Network Components: Weights,
Checkpoints, and {K/V} Caches in Low-Precision Formats''} (2025).

\textbf{General float codecs.} ALP (Afroozeh, Kuffo, and Boncz 2024) and
Pcodec (Loncaric 2025) are strong on smooth or decimal-origin numeric
sequences but target a different data shape than the full-entropy
mantissa of trained weights.

\section{Limitations}\label{limitations}

z4ai cannot exceed the lossless entropy ceiling on dense weights, and
does not try to; its wins are conditional on redundancy being present.
Its pure-Python core compresses several times slower than compiled-C
codecs. As a per-file codec it does not capture cross-model redundancy
at corpus scale, where a dedicated system such as ZipLLM is the better
tool. The high-ratio \texttt{effort="max"} tier trades substantial
decode throughput for a few percent of additional ratio and is therefore
opt-in.

\section{Conclusion}\label{conclusion}

The lossless compressibility of dense trained weights is bounded by
physics, and that wall is already met by existing codecs. The
practically valuable lossless wins for model storage and distribution
come instead from structure, reduced precision, sparsity, and
cross-checkpoint redundancy. z4ai packages the transforms that capture
these --- whole-tensor matching, a lossless palette transform, and
single-reference deltas --- into one self-describing, random-access,
never-worse-than-zstd codec usable without any repository-scale
infrastructure. It is open source and every result here is reproducible
from the repository.

\section*{References}\label{references}
\addcontentsline{toc}{section}{References}

\phantomsection\label{refs}
\begin{CSLReferences}{1}{0}
\bibitem[\citeproctext]{ref-alp}
Afroozeh, Azim, Leonardo X. Kuffo, and Peter Boncz. 2024. {``{ALP}:
Adaptive Lossless Floating-Point Compression.''} In \emph{Proceedings of
the ACM on Management of Data (SIGMOD)}.
\url{https://doi.org/10.1145/3626717}.

\bibitem[\citeproctext]{ref-brotli}
Alakuijala, Jyrki, Andrea Farruggia, Paolo Ferragina, Eugene
Kliuchnikov, Robert Obryk, Zoltan Szabadka, and Lode Vandevenne. 2018.
{``Brotli: A General-Purpose Data Compressor.''} \emph{ACM Transactions
on Information Systems} 37 (1): 1--30.
\url{https://doi.org/10.1145/3231935}.

\bibitem[\citeproctext]{ref-unweight}
Cloudflare Research. 2026. {``Unweight: Lossless {MLP} Weight
Compression for {LLM} Inference.''}
\url{https://research.cloudflare.com/papers/unweight-2026.pdf}.

\bibitem[\citeproctext]{ref-zstd}
Collet, Yann, and Murray S. Kucherawy. 2021. {``Zstandard Compression
and the {'application/zstd'} Media Type.''} RFC 8878. Internet
Engineering Task Force. \url{https://doi.org/10.17487/RFC8878}.

\bibitem[\citeproctext]{ref-rans}
Duda, Jarek. 2013. {``Asymmetric Numeral Systems: Entropy Coding
Combining Speed of {H}uffman Coding with Compression Rate of Arithmetic
Coding.''} \emph{arXiv Preprint arXiv:1311.2540}.
\url{https://doi.org/10.48550/arXiv.1311.2540}.

\bibitem[\citeproctext]{ref-ryg_rans}
Giesen, Fabian. 2014. {``Interleaved Entropy Coders.''} \emph{arXiv
Preprint arXiv:1402.3392}.
\url{https://doi.org/10.48550/arXiv.1402.3392}.

\bibitem[\citeproctext]{ref-neuzip}
Hao, Yongchang, Yanshuai Cao, and Lili Mou. 2024. {``{NeuZip}:
Memory-Efficient Training and Inference with Dynamic Compression of
Neural Networks.''} \emph{Advances in Neural Information Processing
Systems (NeurIPS)}. \url{https://arxiv.org/abs/2410.20650}.

\bibitem[\citeproctext]{ref-zipnn}
Hershcovitch, Moshik, Leshem Choshen, Andrew Wood, Ilias Enmon, Peter
Chin, and Danny Harnik. 2024. {``{ZipNN}: Lossless Compression for {AI}
Models.''} \emph{arXiv Preprint arXiv:2411.05239}.
\url{https://doi.org/10.48550/arXiv.2411.05239}.

\bibitem[\citeproctext]{ref-dietgpu}
Johnson, Jeff. 2022. {``{dietgpu}: GPU-Based Lossless Compression for
Numerical Data.''} \url{https://github.com/facebookresearch/dietgpu}.

\bibitem[\citeproctext]{ref-pcodec}
Loncaric, Martin. 2025. {``Pcodec: Better Compression for Numerical
Sequences.''} \emph{arXiv Preprint arXiv:2502.06112}.
\url{https://arxiv.org/abs/2502.06112}.

\bibitem[\citeproctext]{ref-ckptdelta}
{``Lossless Compression of Neural Network Components: Weights,
Checkpoints, and {K/V} Caches in Low-Precision Formats.''} 2025.
\emph{arXiv Preprint arXiv:2508.19263}.
\url{https://arxiv.org/abs/2508.19263}.

\bibitem[\citeproctext]{ref-zipllm}
Wang et al. 2026. {``{ZipLLM}: Efficient Storage of Large Language Model
Repositories Through Deduplication and Delta Compression.''}
\emph{USENIX Symposium on Networked Systems Design and Implementation
(NSDI)}. \url{https://arxiv.org/abs/2505.06252}.

\bibitem[\citeproctext]{ref-dfloat11}
Zhang, Tianyi, Yang Sui, Shaochen Zhong, Vipin Chaudhary, Xia Hu, and
Anshumali Shrivastava. 2025. {``{DFloat11}: Lossless {LLM} Compression
for Efficient {GPU} Inference via Dynamic-Length Float.''}
\emph{Advances in Neural Information Processing Systems (NeurIPS)}.
\url{https://arxiv.org/abs/2504.11651}.

\bibitem[\citeproctext]{ref-lz77}
Ziv, Jacob, and Abraham Lempel. 1977. {``A Universal Algorithm for
Sequential Data Compression.''} \emph{IEEE Transactions on Information
Theory} 23 (3): 337--43. \url{https://doi.org/10.1109/TIT.1977.1055714}.

\end{CSLReferences}

\end{document}
